\chapter{置換と図形}\label{章:置換の背景}

\chapterref{章:置換積分}では,\,$x=\sin\theta$や$x=\tan\theta$と置いて,根号や分母を簡単にした.しかし,不定積分の答えを$x$へ戻そうとすると,「正弦が$x$になる角」を関数として書く必要がある.また,\,$\sin x$と$\cos x$が分母で混ざった式では,一方を新しい変数にするだけでは,他方が根号として残ることがある.

この章では,まず逆三角関数を定義して,角へ戻る操作を明確にする.続いて,三角関数を同時に有理式へ変える置換を考える.その式は円上の点を直線の傾きで表す式でもあり,整数の直角三角形とも結び付く.最後に,平方の差が一定になる双曲線関数を導入し,同じ根号積分を3つの置換で解く.

1つの積分に解き方が1つとは限らない.この章では別解を並べるだけでなく,2つの解き方の変数の間にどんな置換が成り立つのかを式で確かめ,別解どうしを結ぶ.使う前提は,これまでの三角関数・指数関数の性質と微分,置換積分,部分積分である.新しい関数の微分や積分公式は,それらから導く.

\section{逆三角関数}\label{節:逆三角関数}

$\sin\theta=\frac12$を満たす角は,\,$\frac\pi6$,\,$\frac{5\pi}6$など無数にある.このままでは,\,$x$から角を返す関数にはならない.\sectionref{節:逆向きの置換}で扱った枝の選択を,ここでも行う.

\begin{definition}{逆三角関数とその値の範囲}
次の区間に制限した三角関数の逆関数を,それぞれ$\arcsin$,\,$\arccos$,\,$\arctan$と書く.
\[
\begin{array}{c|c|c}
\text{逆関数}&\text{入力の範囲}&\text{返す角の範囲}\\ \hline
\arcsin x&[-1,1]&\left[-\frac\pi2,\frac\pi2\right]\\[3pt]
\arccos x&[-1,1]&[0,\pi]\\[3pt]
\arctan x&\R&\left(-\frac\pi2,\frac\pi2\right)
\end{array}
\]
例えば$\theta=\arcsin x$とは,\,$\sin\theta=x$かつ$-\frac\pi2\leqq\theta\leqq\frac\pi2$であることを表す.
\end{definition}

正弦と正接は指定した区間で狭義単調増加し,余弦は狭義単調減少するので,角は一つに決まる.正接の範囲の端を含めないのは,そこで$\cos\theta=0$になり,正接が定義されないためである.\,$\sin^{-1}x$などの表記もあるが,逆数との混同を避け,本書では$\arcsin x$を使う.

定義から$\sin(\arcsin x)=x$は$-1\leqq x\leqq1$で成り立つ.一方,\,$\arcsin(\sin\theta)=\theta$は,\,$\theta$が選んだ範囲にあるときに限る.左辺が返す値は,必ずその範囲内だからである.

\begin{example}{逆三角関数の範囲}
次の値を求めよ.
\[
\arcsin\left(\sin\frac{5\pi}6\right),\qquad
\arccos\left(\cos\frac{7\pi}6\right),\qquad
\arctan\left(\tan\frac{3\pi}4\right).
\]
\end{example}

第一の正弦は$\frac12$であり,\,$\arcsin$の値の範囲にある角は$\frac\pi6$である.第二の余弦は$-\frac{\sqrt3}2$であり,\,$[0,\pi]$の中では$\frac{5\pi}6$になる.第三の正接は$-1$であり,\,$\left(-\frac\pi2,\frac\pi2\right)$の中では$-\frac\pi4$になる.従って答えは順に
\[
\frac\pi6,\qquad\frac{5\pi}6,\qquad-\frac\pi4
\]
である.逆関数の記号を打ち消す前に,戻った角が指定された範囲にあるかを確かめる.

\subsection{逆三角関数の微分}

逆関数の微分は,もとの三角関数の微分から求められる.\,$y=\arcsin x$,\,$-1<x<1$とおく.\,$\sin y=x$の両辺を$x$で微分すると
\[
\cos y\cdot y'=1
\]
となる.ここで$-\frac\pi2<y<\frac\pi2$なので$\cos y>0$である.したがって$\cos y=\sqrt{1-\sin^2y}=\sqrt{1-x^2}$と符号まで決まり,
\[
(\arcsin x)'=\frac1{\sqrt{1-x^2}}\qquad(-1<x<1)
\]
を得る.端点では分母が$0$になり,この式は使えない.

$y=\arccos x$の場合には,\,$\cos y=x$を微分して$-\sin y\cdot y'=1$となる.\,$0<y<\pi$で$\sin y=\sqrt{1-x^2}>0$なので
\[
(\arccos x)'=-\frac1{\sqrt{1-x^2}}\qquad(-1<x<1)
\]
である.負号は,余弦が選んだ区間で減少することとも一致する.\,$y=\arctan x$では$\tan y=x$を微分して
\[
(1+\tan^2y)y'=1
\]
となり,\,$\tan y=x$を代入して
\[
(\arctan x)'=\frac1{1+x^2}\qquad(x\in\R)
\]
を得る.

合成関数の微分を使えば,\,$a>0$に対して
\begin{equation}\label{式:逆三角関数の積分}
\begin{aligned}
\int\frac1{a^2+x^2}\dd x
 &=\frac1a\arctan\frac xa+C,\\
\int\frac1{\sqrt{a^2-x^2}}\dd x
 &=\arcsin\frac xa+C\qquad(|x|<a)
\end{aligned}
\end{equation}
を得る.第1式では,\,$x$を$a$で割ることと,外側の係数$\frac1a$の両方が必要である.第2式では,根号から出る$a$と合成関数の微分の$\frac1a$が打ち消し合う.

一般の2次式が分母や根号の中にあるときは,平方完成して$(x-p)^2+a^2$や$a^2-(x-p)^2$の形を作り,\,$u=\frac{x-p}{a}$と置けば,上の公式の$a=1$の場合に帰着する.

\begin{bookpoint}
\textbf{POINT}\par
2次式を$(x-p)^2+a^2$や$a^2-(x-p)^2$と平方完成し,\,$u=\frac{x-p}{a}$と置いて$\arctan u$や$\arcsin u$の微分の形にする.
\end{bookpoint}

\begin{example}{二次式の逆数の積分}
次の不定積分を求めよ.
\[
\text{(1)}\quad\int\frac1{x^2-4x+8}\dd x,
\qquad
\text{(2)}\quad\int\frac1{\sqrt{9-4x^2}}\dd x.
\]
\end{example}

\paragraphbf{方針}
\begin{enumerate}[label=(\arabic*)]
\item 分母を$(x-2)^2+4$と平方完成し,\,$u=\frac{x-2}{2}$と置く.
\item 根号の中を$9\left\{1-\left(\frac{2x}{3}\right)^2\right\}$と書き,\,$u=\frac{2x}{3}$と置く.
\end{enumerate}

\paragraphbf{解答}
\begin{enumerate}[label=(\arabic*)]
\item $\frac12\arctan\frac{x-2}{2}+C$
\item $\frac12\arcsin\frac{2x}{3}+C$\quad$\left(-\frac32<x<\frac32\right)$
\end{enumerate}

\paragraphbf{解法}
\begin{enumerate}[label=(\arabic*)]
\item 分母は$x^2-4x+8=(x-2)^2+4$である.\,$u=\frac{x-2}{2}$と置くと,\,$(x-2)^2+4=4(1+u^2)$,\,$\dd x=2\dd u$だから
\begin{align*}
\int\frac1{x^2-4x+8}\dd x
&=\int\frac{2}{4(1+u^2)}\dd u\\
&=\frac12\arctan u+C\\
&=\frac12\arctan\frac{x-2}{2}+C.
\end{align*}
\item もとの関数が実数で定義されるのは$-\frac32<x<\frac32$である.\,$u=\frac{2x}{3}$と置くと,\,$\sqrt{9-4x^2}=3\sqrt{1-u^2}$,\,$\dd x=\frac32\dd u$だから
\begin{align*}
\int\frac1{\sqrt{9-4x^2}}\dd x
&=\int\frac{3}{2\cdot3\sqrt{1-u^2}}\dd u\\
&=\frac12\arcsin u+C\\
&=\frac12\arcsin\frac{2x}{3}+C.
\end{align*}
\end{enumerate}

\paragraphbf{三角置換との関係}
(1)は,\,\chapterref{章:置換積分}の三角置換でも解ける.\,$x-2=2\tan\theta$,\,$-\frac\pi2<\theta<\frac\pi2$と置くと,\,$\dd x=2(1+\tan^2\theta)\dd\theta$,\,$(x-2)^2+4=4(1+\tan^2\theta)$だから
\[
\int\frac1{x^2-4x+8}\dd x
=\int\frac12\dd\theta
=\frac\theta2+C
\]
となる.ここで$\theta$は$\tan\theta=\frac{x-2}{2}$を満たす$\left(-\frac\pi2,\frac\pi2\right)$の角,すなわち$\theta=\arctan\frac{x-2}{2}$である.2つの解き方は,\,$u=\tan\theta$という置換で結ばれている.逆三角関数は,三角置換で導入した角$\theta$を$x$の関数として書くための記号なのだ.
%
%
\section{三角関数の逆数}\label{節:相反三角関数}

三角置換では$\frac1{\cos\theta}$や$\frac1{\sin\theta}$が繰り返し現れる.これらを一つの関数として書くと,平方や微分の関係を読み取りやすい.
\begin{definition}{正割・余割・余接}
分母が$0$でないところで
\[
\sec x=\frac1{\cos x},\qquad
\csc x=\frac1{\sin x},\qquad
\cot x=\frac{\cos x}{\sin x}
\]
と定義する.順に正割,余割,余接という.ここでの逆数と,前節の逆関数は異なる操作である.
\end{definition}

$\sin^2x+\cos^2x=1$を$\cos^2x$,\,$\sin^2x$で割ると,それぞれ
\[
1+\tan^2x=\sec^2x,\qquad
1+\cot^2x=\csc^2x
\]
となる.微分も,すでに知っている商の微分から得られる.例えば
\[
(\sec x)'=\frac{\sin x}{\cos^2x}=\sec x\tan x.
\]
同様に計算して
\begin{equation}\label{式:相反三角関数の微分}
(\csc x)'=-\csc x\cot x,\qquad
(\cot x)'=-\csc^2x
\end{equation}
を得る.この記法により,\,$x=\tan\theta$と置いたときの$\dd x$は$\sec^2\theta\dd\theta$,\,$\sqrt{1+x^2}$は$|\sec\theta|$と短く書ける.ただし,記号を短くしても,絶対値を外すための符号の確認は残る.

\begin{example}{三角関数の逆数の積分}
$\int\sec x\dd x$と$\int\csc x\dd x$を求めよ.
\end{example}

$\sec x$の微分は$\sec x\tan x$,\,$\tan x$の微分は$\sec^2x$である.この二つを足すと
\[
(\sec x+\tan x)'=\sec x(\sec x+\tan x)
\]
となり,元の$\sec x$が因子に現れる.また
\[
(\sec x+\tan x)(\sec x-\tan x)=1
\]
だから,\,$\sec x+\tan x$は定義されるところで$0$にならない.従って
\[
\int\sec x\dd x
=\int\frac{(\sec x+\tan x)'}{\sec x+\tan x}\dd x
=\log|\sec x+\tan x|+C.
\]
これは,分母の微分を分子に作る対数型の積分である.

余割では引き算を考えると
\[
(\csc x-\cot x)'=-\csc x\cot x+\csc^2x
=\csc x(\csc x-\cot x).
\]
さらに$(\csc x-\cot x)(\csc x+\cot x)=1$なので,
\[
\int\csc x\dd x=\log|\csc x-\cot x|+C
\]
となる.どちらも,もとの三角関数が定義される一つの区間で使う.特異点をまたぐ不定積分の定数を,一つにそろえる必要はない.


%
\section{半角置換}\label{節:半角置換}

正弦と余弦を足す・掛ける・割ることで作った式の積分を考える.このような式は,2変数の\textbf{有理式}$R(X,Y)$の$X,Y$に$\sin x,\cos x$を代入したものなので,一般的な形は
\[
\int R(\sin x,\cos x)\dd x
\]
で与えられる.これを,\,$t$の有理式の積分に帰着させることを目標とする.有理式なら,割り算や部分分数分解で積分できるからである.

さて,何と置けばよいだろうか.候補を挙げてみよう.
\begin{enumerate}
\item $t=\sin x$と置くと,\,$\cos x=\pm\sqrt{1-t^2}$となり,根号が残る.
\item $t=\cos x$と置いても,同じく\,$\sin x$が根号で残る.
\item $t=\tan x$と置くと,\,$\cos^2x=\frac1{1+t^2}$,\,$\sin x\cos x=\frac t{1+t^2}$は有理式になる.しかし$\sin x$単独は$\pm\frac{t}{\sqrt{1+t^2}}$となり,根号が残る.
\end{enumerate}
3では,\,$\cos^2x$と$\sin x\cos x$は有理式になった.したがって,2倍角の公式から
\[
\sin2x=2\sin x\cos x=\frac{2t}{1+t^2},\qquad
\cos2x=2\cos^2x-1=\frac{1-t^2}{1+t^2}
\]
も$t=\tan x$の有理式である.つまり,正接は2倍の角の正弦・余弦を有理式で表す.そこで角を半分にして,\,$\tan\frac x2$を使えばよい.

$t=\tan\frac x2$と置く.2倍角の公式で分子を書き,分母の$1=\cos^2\frac x2+\sin^2\frac x2$で割ると
\begin{equation}\label{式:半角置換}
\begin{aligned}
\sin x&=\cfrac{2\sin\cfrac x2\cos\cfrac x2}
 {\cos^2\cfrac x2+\sin^2\cfrac x2}
 =\frac{2t}{1+t^2},\\
\cos x&=\cfrac{\cos^2\cfrac x2-\sin^2\cfrac x2}
 {\cos^2\cfrac x2+\sin^2\cfrac x2}
 =\frac{1-t^2}{1+t^2},\\
\frac{\dd t}{\dd x}&=\frac12\sec^2\frac x2
 =\frac{1+t^2}{2},\qquad
\dd x=\frac{2}{1+t^2}\dd t
\end{aligned}
\end{equation}
を得る.最初の2式では,中央の分数の分子と分母を$\cos^2\frac x2$で割った.

ここで1つ断りをいれておこう.この割り算と$\tan\frac x2$の定義には,\,$\cos\frac x2\ne0$,すなわち$x\ne(2k+1)\pi$($k$は整数)が必要である.そこで,この置換は$(2k-1)\pi<x<(2k+1)\pi$の1つの区間の中で使う.この区間では$t$は実数全体を狭義単調増加して動き,逆向きには
\[
x=2\arctan t+2k\pi
\]
と書ける.とくに$-\pi<x<\pi$では$x=2\arctan t$である.

この置換を\textbf{半角置換},または\textbf{ワイエルシュトラス置換}という.正弦・余弦・$\dd x$のすべてが$t$の有理式になるので,置換後は有理式の積分になる.

\begin{bookpoint}
\textbf{POINT}\par
$t=\tan\frac x2$と置き,\,$\sin x=\frac{2t}{1+t^2}$,\,$\cos x=\frac{1-t^2}{1+t^2}$,\,$\dd x=\frac{2}{1+t^2}\dd t$を代入して$t$の有理式にする.
\end{bookpoint}

\begin{example}{半角置換}
$\int_0^{\frac\pi2}\frac1{2+\sin x}\dd x$を求めよ.
\end{example}

\paragraphbf{方針}
分母$2+\sin x$の微分$\cos x$は分子に現れないので,\,$t=\tan\frac x2$と置いて$t$の有理式にする.最後は平方完成して$\arctan$で積分する.

\paragraphbf{解答}
$\frac\pi{3\sqrt3}$

\paragraphbf{解法}
$t=\tan\frac x2$と置くと,\,$x=0,\frac\pi2$は$t=0,1$へ移る.式\eqref{式:半角置換}を代入すると
\begin{align*}
\int_0^{\frac\pi2}\frac1{2+\sin x}\dd x
&=\int_0^1\cfrac1{2+\cfrac{2t}{1+t^2}}\cdot
 \frac2{1+t^2}\dd t\\
&=\int_0^1\frac{2}{2(1+t^2)+2t}\dd t\\
&=\int_0^1\frac1{t^2+t+1}\dd t
\end{align*}
となる.分母を平方完成すると$t^2+t+1=\left(t+\frac12\right)^2+\frac34$なので
\begin{align*}
\int_0^1\frac1{t^2+t+1}\dd t
&=\frac2{\sqrt3}\left[\arctan\frac{2t+1}{\sqrt3}\right]_0^1\\
&=\frac2{\sqrt3}\left(\frac\pi3-\frac\pi6\right)\\
&=\frac\pi{3\sqrt3}
\end{align*}
を得る.

\paragraphbf{区間をまたぐ場合}
もし積分区間が$[0,2\pi]$なら,\,$x=\pi$で$t$が定義されなくなる.端点だけを$t=\tan0=0$,\,$t=\tan\pi=0$へ換えて答えを$0$としてはいけない.元の被積分関数は$\pi$で連続でも,置換に使った$t$はそこで途切れる.この場合には$\pi$の前後に分け,\,$t\to+\infty$と$t\to-\infty$を別々に扱う.章末問題で確かめる.

\paragraphbf{正割の積分の別解}
前節の$\int\sec x\dd x$も,半角置換で解ける.\,$-\pi<x<\pi$,\,$x\ne\pm\frac\pi2$で$t=\tan\frac x2$と置くと,\,$\sec x=\frac{1+t^2}{1-t^2}$なので
\begin{align*}
\int\sec x\dd x
&=\int\frac{1+t^2}{1-t^2}\cdot\frac{2}{1+t^2}\dd t\\
&=\int\left(\frac1{1+t}+\frac1{1-t}\right)\dd t\\
&=\log\left|\frac{1+t}{1-t}\right|+C
\end{align*}
となる.一方,\,$\sec x+\tan x$を$t$で書くと
\begin{align*}
\sec x+\tan x
&=\frac{1+\sin x}{\cos x}\\
&=\cfrac{1+\cfrac{2t}{1+t^2}}{\cfrac{1-t^2}{1+t^2}}\\
&=\frac{(1+t)^2}{(1+t)(1-t)}
=\frac{1+t}{1-t}
\end{align*}
である.前節で結果からお見せした$\sec x+\tan x$は,半角置換で書けば$\frac{1+t}{1-t}$という簡単な有理式なのだ.
%
\section{円の媒介変数表示}\label{節:円の有理パラメータ}

前節の2式は
\[
X=\frac{1-t^2}{1+t^2},\qquad
Y=\frac{2t}{1+t^2}
\]
と書くと,\,$X^2+Y^2=1$を満たす.正弦と余弦を使ったのだから当然ともいえるが,この式は三角関数なしでも導ける.

図\ref{図:円の有理パラメータ}のように,単位円上の点$A=(-1,0)$を固定し,傾き$t$の直線$Y=t(X+1)$を引く.円との交点を求めるため代入すると
\[
\begin{aligned}
X^2+t^2(X+1)^2&=1,\\
(X+1)\{(1+t^2)X+(t^2-1)\}&=0.
\end{aligned}
\]
1つの解$X=-1$は,固定した点$A$である.もう1つの交点$P$は,上の$X,Y$で表される.\,$t$を動かして点を表すので,\,$t$を\textbf{媒介変数},この表示を円の\textbf{有理パラメータ表示}という.

\begin{figure}[htbp]
\centering
\begin{tikzpicture}[x=2.25cm,y=2.25cm,>=Stealth]
 \draw[->] (-1.4,0)--(1.45,0) node[right]{$X$};
 \draw[->] (0,-1.2)--(0,1.35) node[above]{$Y$};
 \draw[thick,tealblack] (0,0) circle(1);
 \draw[thick,blueblack] (-1.23,-0.115)--(1.15,1.075);
 \draw[dashed] (0.6,0)--(0.6,0.8)--(0,0.8);
 \draw[blueblack] (0,0)--(0.6,0.8);
 \draw[->] (0.26,0) arc[start angle=0,end angle=53.1301,radius=0.26];
 \node at (0.35,0.17) {$x$};
 \draw[->] (-0.65,0) arc[start angle=0,end angle=26.5651,radius=0.35];
 \node at (-0.42,0.15) {$\frac x2$};
 \fill (-1,0) circle(1.2pt) node[below left]{$A=(-1,0)$};
 \fill (0.6,0.8) circle(1.2pt);
 \draw[thin,gray] (0.63,0.79)--(1.04,0.56);
 \node[anchor=west] at (1.04,0.52) {$P=(\cos x,\sin x)$};
 \node[above right] at (0.97,0.985) {$Y=t(X+1)$};
 \node[below left] at (0,0) {$O$};
\end{tikzpicture}
\caption{円上の固定点から引いた直線の傾きで,もう1つの交点を表す.図は$t=\frac12$の場合.}
\label{図:円の有理パラメータ}
\end{figure}

$P=(X,Y)\ne A$からは$t=\frac{Y}{X+1}$と戻せるので,有限な実数$t$で,円上の$A$以外のすべての点を一度ずつ表せる.\,$P=(\cos x,\sin x)$なら
\[
t=\frac{\sin x}{1+\cos x}=\tan\frac x2
\]
である.最後の等式は2倍角の公式から確かめられる.図の上半円では,三角形$AOP$の角を見ても,直線$AP$が正の$X$方向となす角は$\frac x2$になる.実際,\,$OA=OP=1$より三角形$AOP$は二等辺三角形であり,頂角が$\angle AOP=\pi-x$なので,底角$\angle OAP$は$\frac{\pi-(\pi-x)}{2}=\frac x2$である.\,$A$から$O$へ向かう向きは正の$X$方向なので,この底角が求める角である.半角が現れる理由を,円の上の角としても読めた.表せなかった点$A$は,ちょうど$x=(2k+1)\pi$に対応する.

\subsection{ピタゴラス数}

$t=\frac nm$を整数$m,n$で表すと,
\[
X=\frac{m^2-n^2}{m^2+n^2},\qquad
Y=\frac{2mn}{m^2+n^2}.
\]
従って
\begin{equation}\label{式:ピタゴラス数}
(m^2-n^2)^2+(2mn)^2=(m^2+n^2)^2
\end{equation}
である.\,$m>n>0$なら三つの整数は正であり,直角三角形の三辺になる.\,$a^2+b^2=c^2$を満たす正の整数の組を\textbf{ピタゴラス数}という.

逆に整数の直角三角形$(a,b,c)$があれば,点$\left(\frac ac,\frac bc\right)$は単位円上にある.その傾きは
\[
t=\cfrac{\cfrac bc}{1+\cfrac ac}=\frac{b}{c+a}
\]
と有理数になる.これを$\frac nm$と表せば,上の公式が元の三辺と同じ比を与える.従って,この表示は偶然いくつかの例を作るだけでなく,すべての整数の直角三角形の辺の比を表す.得られた三整数に共通因数があれば,それで割ればよい.

三辺の最大公約数が$1$のピタゴラス数を,\textbf{原始ピタゴラス数}という.整数$m>n>0$が互いに素で,一方が偶数,他方が奇数ならば,式\eqref{式:ピタゴラス数}は原始ピタゴラス数を与える.逆に,すべての原始ピタゴラス数は,直角を挟む二辺の交換を除いてこの形で表せる.この主張を,整数の素因数分解の一意性を使って確かめよう.

まず上の条件で作った三整数に,共通の素因数$p$があると仮定する.\,$m^2-n^2$と$m^2+n^2$はともに奇数なので,\,$p$は奇素数である.この二数の和と差から,\,$p$は$2m^2$と$2n^2$を割る.従って$p$は$m,n$の両方を割り,互いに素という条件に反する.共通の素因数がないので,三整数の最大公約数は$1$である.

逆に原始ピタゴラス数$(a,b,c)$を考える.二辺$a,b$がともに偶数なら原始ではなく,ともに奇数なら$a^2+b^2$を$4$で割った余りが$2$となり,平方数$c^2$にはなれない.従って一方だけが偶数であり,必要なら辺を交換して$b$を偶数,\,$a$を奇数とする.このとき$c$も奇数である.また$a,c$に共通の素因数があれば,\,$b^2=c^2-a^2$,従って$b$もその素因数で割れ,原始という条件に反する.よって$a,c$は互いに素である.

ここで
\[
A=\frac{c+a}{2},\qquad B=\frac{c-a}{2}
\]
と置くと,\,$A>B>0$の整数であり,
\[
AB=\frac{c^2-a^2}{4}=\left(\frac b2\right)^2.
\]
$A,B$の共通因数は$A+B=c$と$A-B=a$も割るので,\,$A,B$は互いに素である.互いに素な二数の積が平方数なら,素因数は片方にしか現れず,その指数はすべて偶数になる.従って$A=m^2$,\,$B=n^2$と書ける.和と差を戻せば
\[
a=m^2-n^2,\qquad b=2mn,\qquad c=m^2+n^2.
\]
$A,B$が互いに素だから$m,n$も互いに素であり,\,$c$が奇数だから$m,n$の偶奇は異なる.これで逆向きも示せた.円の表示が辺の比を与えることに加えて,約分の要らない整数を作る条件まで分かった.

\begin{example}{ピタゴラス数}
$t=\frac23$からピタゴラス数を作れ.また三辺が$8,15,17$の三角形に対応する傾きを求めよ.
\end{example}

$m=3,n=2$なので
\[
m^2-n^2=5,\qquad2mn=12,\qquad m^2+n^2=13.
\]
確かに$5^2+12^2=13^2$である.後半では$P=\left(\frac{15}{17},\frac8{17}\right)$と選ぶと
\[
t=\frac8{17+15}=\frac14.
\]
$m=4,n=1$を公式へ戻せば$(15,8,17)$となる.二つの直角を挟む辺のどちらを$X$座標にするかで,対応する傾きは変わる.



%
%
\section{双曲線関数}\label{節:双曲線関数}

円では$X^2+Y^2=1$を使った.これに対して,\,$\sqrt{1+x^2}$を簡単にしたい場合には,\,$Y^2-X^2=1$という平方の差を使えればよい.指数関数の和と差が,その形を作る.

\begin{definition}{双曲線関数}
実数$u$に対して
\[
\sinh u=\frac{\e^u-\e^{-u}}2,\qquad
\cosh u=\frac{\e^u+\e^{-u}}2,\qquad
\tanh u=\frac{\sinh u}{\cosh u}
\]
と定義する.順に双曲線正弦,双曲線余弦,双曲線正接という.
\end{definition}

$\cosh u>0$なので$\tanh u$はすべての実数で定義される.和と差の積から
\[
\begin{aligned}
\cosh^2u-\sinh^2u
&=\frac{(\e^u+\e^{-u})^2-(\e^u-\e^{-u})^2}{4}\\
&=\e^u\e^{-u}=1
\end{aligned}
\]
となる.したがって$(X,Y)=(\sinh u,\cosh u)$は,双曲線$Y^2-X^2=1$の上側にある.図\ref{図:双曲線の媒介表示}のように,\,$u$が正なら$X>0$,負なら$X<0$である.この$u$は,円のときのような通常の中心角として定義したものではない.

\begin{figure}[htbp]
\centering
\begin{tikzpicture}[x=1.5cm,y=1.5cm,>=Stealth]
 \draw[->] (-2.4,0)--(2.55,0) node[right]{$X$};
 \draw[->] (0,-0.15)--(0,2.8) node[above]{$Y$};
 \draw[dashed,gray] (-2.2,2.2)--(0,0)--(2.2,2.2);
 \draw[very thick,tealblack,domain=-2.2:2.2,samples=80]
  plot(\x,{sqrt(1+\x*\x)});
 \draw[dashed,blueblack] (1.175201,0)--(1.175201,1.543081)--(0,1.543081);
 \fill[blueblack] (1.175201,1.543081) circle(1.3pt);
 \node[anchor=west] at (0.25,2.55) {$Y^2-X^2=1$};
 \node[anchor=east] at (-0.3,1.65) {$u<0$};
 \node[anchor=west] at (0.35,1.95) {$u>0$};
 \node[anchor=west] at (1.4,0.58) {$(\sinh u,\cosh u)$};
 \draw[->,blueblack] (1.75,0.82)--(1.2,1.48);
 \node[left] at (0,1) {$1$};
 \node[below left] at (0,0) {$0$};
\end{tikzpicture}
\caption{指数関数の和と差によって,双曲線の上側を一度ずつ表す.点は$u=1$の位置.}
\label{図:双曲線の媒介表示}
\end{figure}

定義を微分すると
\[
(\sinh u)'=\frac{\e^u+\e^{-u}}2=\cosh u,
\qquad
(\cosh u)'=\frac{\e^u-\e^{-u}}2=\sinh u.
\]
余弦の場合と違って,\,$\cosh$を微分しても負号は付かない.商の微分から
\begin{equation}\label{式:双曲線正接の微分}
(\tanh u)'=\frac{\cosh^2u-\sinh^2u}{\cosh^2u}
=\frac1{\cosh^2u}=1-\tanh^2u
\end{equation}
である.加えて指数関数の定義式を展開すると
\begin{equation}\label{式:双曲線の二倍公式}
\sinh(2u)=2\sinh u\cosh u,\qquad
\cosh(2u)=2\cosh^2u-1
\end{equation}
となる.例えば後者の右辺は
\[
\frac{(\e^u+\e^{-u})^2}{2}-1
=\frac{\e^{2u}+\e^{-2u}}2
\]
であり,左辺の定義に一致する.

\subsection{双曲線関数の逆関数}

$(\sinh u)'=\cosh u>0$なので,\,$\sinh u$は実数全体で狭義単調増加する.また$u\to\infty$で$\sinh u\to\infty$,\,$u\to-\infty$で$\sinh u\to-\infty$だから,どの実数$x$も$x=\sinh u$と一意に書ける.この逆関数を$\operatorname{arsinh}x$と書く.

$x=\sinh u$なら,\,$\cosh u>0$と$\cosh^2u=1+\sinh^2u$より$\cosh u=\sqrt{1+x^2}$である.定義式を足すと$\sinh u+\cosh u=\e^u$だから
\begin{equation}\label{式:逆双曲線正弦}
u=\operatorname{arsinh}x=\log\bigl(x+\sqrt{1+x^2}\bigr)
\end{equation}
となる.対数の中は,\,$\sqrt{1+x^2}>|x|$によって常に正である.逆関数の微分から
\[
\frac{\dd u}{\dd x}=\frac1{\cosh u}
=\frac1{\sqrt{1+x^2}}
\]
を得る.ここでも積分の値から逆関数を定義したのではなく,指数関数による定義から順に導いた.

この関係を置換積分に使おう.\,$a>0$として$x=a\sinh u$と置くと
\[
\sqrt{a^2+x^2}=a\cosh u,\qquad
\dd x=a\cosh u\dd u
\]
となり,根号と$\dd x$に同じ$\cosh u$が現れる.三角置換$x=a\tan\theta$では,根号が$a\sec\theta$,\,$\dd x$が$a\sec^2\theta\dd\theta$となって次数がずれたが,双曲線関数では次数がずれない.

\begin{bookpoint}
\textbf{POINT}\par
$\sqrt{a^2+x^2}$を含む積分は$x=a\sinh u$と置き,\,$\sqrt{a^2+x^2}=a\cosh u$,\,$\dd x=a\cosh u\dd u$を代入する.
\end{bookpoint}

\begin{example}{双曲線正弦への置換}
$\int_0^3\frac1{\sqrt{4+x^2}}\dd x$を求めよ.
\end{example}

\paragraphbf{方針}
$x=2\sinh u$と置くと,分母と$\dd x$の$2\cosh u$が打ち消し合い,\,$1$の積分になる.

\paragraphbf{解答}
$\log\frac{3+\sqrt{13}}2$

\paragraphbf{解法}
$x=2\sinh u$と置けば
\[
\dd x=2\cosh u\dd u,\qquad
\sqrt{4+x^2}=2\cosh u
\]
である.端点$x=0$は$u=0$へ,\,$x=3$は
\[
u=\operatorname{arsinh}\frac32
=\log\left(\frac32+\sqrt{1+\frac94}\right)
=\log\frac{3+\sqrt{13}}2
\]
へ移る.したがって
\[
\int_0^3\frac1{\sqrt{4+x^2}}\dd x
=\int_0^{\log\frac{3+\sqrt{13}}2}1\dd u
=\log\frac{3+\sqrt{13}}2
\]
となる.

\paragraphbf{三角置換との関係}
同じ積分を$x=2\tan\theta$で置換すると,\,$\int\sec\theta\dd\theta$が残る.\,\sectionref{節:相反三角関数}の結果を使えば
\[
\int\sec\theta\dd\theta=\log(\sec\theta+\tan\theta)+C
\]
である.ここで$\tan\theta=\frac x2$,\,$\sec\theta=\frac{\sqrt{4+x^2}}2$を代入すると,対数の中は$\frac{x+\sqrt{4+x^2}}2$となり,\,$\operatorname{arsinh}\frac x2$と一致する.2つの置換は$\tan\theta=\sinh u$で結ばれ,このとき$\sec\theta=\cosh u$,\,$\sec\theta+\tan\theta=\e^u$となる.正割の積分に現れた$\sec\theta+\tan\theta$は,双曲線関数で書けば$\e^u$だったのだ.

また\chapterref{章:部分積分}では,\,$\log\bigl(x+\sqrt{1+x^2}\bigr)$を微分すると$\frac1{\sqrt{1+x^2}}$になることを直接確かめた.式\eqref{式:逆双曲線正弦}は,その対数が$\operatorname{arsinh}x$であることを表している.

\paragraphbf{偶関数の逆関数}
$\cosh u$は偶関数なので,実数全体では逆関数を持たない.しかし$u\geqq0$に制限すれば,\,$1$から無限大へ狭義単調増加する.\,$x=\cosh u$を使う置換では,このように枝を選ぶ.\,$\sinh$と$\cosh$で,逆向きに戻すときの条件は同じではない.

\begin{example}{双曲線正接への置換}
$\int_0^{\frac12}\frac1{1-x^2}\dd x$を求めよ.
\end{example}

\paragraphbf{方針}
$x=\tanh u$と置くと,\,$\dd x=(1-x^2)\dd u$となって被積分関数の分母が消える.

\paragraphbf{解答}
$\frac12\log3$

\paragraphbf{解法}
$\tanh u=\frac{\e^{2u}-1}{\e^{2u}+1}$は実数全体で狭義単調増加し,値の範囲は$(-1,1)$である.\,$x=\tanh u$と置けば,式\eqref{式:双曲線正接の微分}によって
\[
\dd x=(1-\tanh^2u)\dd u=(1-x^2)\dd u
\]
となる.元へ戻す式は,\,$x=\frac{\e^{2u}-1}{\e^{2u}+1}$を$\e^{2u}$について解いて
\begin{align*}
x(\e^{2u}+1)&=\e^{2u}-1\\
\e^{2u}&=\frac{1+x}{1-x}\\
u&=\frac12\log\frac{1+x}{1-x}
\end{align*}
である.したがって端点$x=0,\frac12$は$u=0,\frac12\log3$へ移り,
\[
\int_0^{\frac12}\frac1{1-x^2}\dd x
=\int_0^{\frac12\log3}1\dd u
=\frac12\log3
\]
となる.

\paragraphbf{部分分数分解との関係}
部分分数分解で
\[
\frac1{1-x^2}=\frac12\left(\frac1{1+x}+\frac1{1-x}\right)
\]
と分けても解ける.原始関数は
\[
\frac12\log(1+x)-\frac12\log(1-x)
=\frac12\log\frac{1+x}{1-x}
\]
であり,\,$-1<x<1$で上の$u$と同じ式である.部分分数分解で現れる2つの対数の差が,双曲線正接の逆関数そのものになっている.
%
%
\section{根号を含む積分の別解}\label{節:根号積分の三方法}

ここまでの道具を,同じ積分に使ってみよう.違いは答えそのものではなく,途中のどの式が簡単になるかにある.

\begin{example}{三つの置換による別解}
$\int\sqrt{1+x^2}\dd x$を,三角置換,双曲線関数の置換,根号を含む代数的な置換によって求めよ.
\end{example}

\subsection{三角置換}

$x=\tan\theta$,\,$-\frac\pi2<\theta<\frac\pi2$と置く.この範囲では$\sec\theta>0$なので
\[
\sqrt{1+x^2}=\sec\theta,\qquad
\dd x=\sec^2\theta\dd\theta,\qquad
\int\sqrt{1+x^2}\dd x=\int\sec^3\theta\dd\theta.
\]
根号はなくなったが,三乗の積分が残った.\,$J=\int\sec^3\theta\dd\theta$と書き,\,$\sec\theta$を微分する側,\,$\sec^2\theta$を積分する側にして部分積分する.
\[
\begin{aligned}
J&=\sec\theta\tan\theta
-\int\sec\theta\tan^2\theta\dd\theta\\
&=\sec\theta\tan\theta
-\int\sec\theta(\sec^2\theta-1)\dd\theta\\
&=\sec\theta\tan\theta-J+\int\sec\theta\dd\theta.
\end{aligned}
\]
任意定数を最後にまとめるものとして,この等式を$J$について解く.前に求めた正割の積分から
\[
J=\frac12\sec\theta\tan\theta
+\frac12\log(\sec\theta+\tan\theta)+C.
\]
この範囲では$\sec\theta>|\tan\theta|$だから,対数の中は正である.\,$x$へ戻して
\begin{equation}\label{式:一と二乗の根号積分}
\int\sqrt{1+x^2}\dd x
=\frac12\left\{x\sqrt{1+x^2}
+\log\bigl(x+\sqrt{1+x^2}\bigr)\right\}+C
\end{equation}
を得る.三角置換だけで計算が終わるのではなく,その後に部分積分を使った.


%
\subsection{双曲線置換}

$x=\sinh u$と置くと,実数$x$全体に対して$u$が一つに決まり,
\[
\sqrt{1+x^2}=\cosh u,\qquad
\dd x=\cosh u\dd u
\]
となる.今回は微分も同じ$\cosh u$になるので,積分は二乗で済む.式\eqref{式:双曲線の二倍公式}より
\[
\begin{aligned}
\int\sqrt{1+x^2}\dd x
&=\int\cosh^2u\dd u\\
&=\frac12\int\{1+\cosh(2u)\}\dd u\\
&=\frac u2+\frac14\sinh(2u)+C\\
&=\frac12\{u+\sinh u\cosh u\}+C.
\end{aligned}
\]
$u=\log(x+\sqrt{1+x^2})$と戻せば,式\eqref{式:一と二乗の根号積分}と一致する.


%
\subsection{根号を含む式への置換}

第3の方法では,関数の名前を使わずに
\[
t=x+\sqrt{1+x^2}
\]
と置く.この和を選ぶ理由は,差との積が$1$になるためである.
\[
\bigl(\sqrt{1+x^2}+x\bigr)
\bigl(\sqrt{1+x^2}-x\bigr)=1.
\]
したがって$\sqrt{1+x^2}-x=\frac1t$であり,この式と$t$の式の和と差を取って2で割れば
\[
x=\frac12\left(t-\frac1t\right),\qquad
\sqrt{1+x^2}=\frac12\left(t+\frac1t\right),\qquad
\dd x=\frac12\left(1+\frac1{t^2}\right)\dd t.
\]
$t>0$であり,この$x$の式は$t>0$から実数全体への狭義単調増加な対応である.根号も微分も有理式へ変わったので,
\[
\begin{aligned}
\int\sqrt{1+x^2}\dd x
&=\frac14\int\left(t+\frac2t+\frac1{t^3}\right)\dd t\\
&=\frac{t^2}{8}+\frac12\log t-\frac1{8t^2}+C.
\end{aligned}
\]
ここで,\,$x$と根号の$t$による式を掛けると
\[
x\sqrt{1+x^2}=\frac14\left(t^2-\frac1{t^2}\right)
\]
なので,\,$\frac{t^2}{8}-\frac1{8t^2}=\frac12x\sqrt{1+x^2}$である.また$\frac12\log t=\frac12\log\bigl(x+\sqrt{1+x^2}\bigr)$なので,式\eqref{式:一と二乗の根号積分}と同じ答えになる.

\paragraphbf{3つの置換のつながり}
3つの置換は,変数どうしの置換で直接結ばれている.まず双曲線置換と代数的な置換では,\,$x=\sinh u$なら
\[
t=x+\sqrt{1+x^2}=\sinh u+\cosh u=\e^u
\]
となる.つまり第3の方法は,双曲線の媒介変数$u$の代わりに$t=\e^u$を変数として,\,$\e^u$と$\e^{-u}$の式を有理式に直したものである.実際,\,$\cosh^2u=\frac14\left(t+\frac1t\right)^2$,\,$\dd u=\frac{\dd t}{t}$を代入すれば,第2の方法の$\int\cosh^2u\dd u$は第3の方法の$t$の積分になる.

次に三角置換と双曲線置換では,\,$x=\tan\theta=\sinh u$であり,
\[
\sec\theta=\sqrt{1+\tan^2\theta}=\sqrt{1+\sinh^2u}=\cosh u
\]
となる.\,$\dd x$を2通りに書くと$\sec^2\theta\dd\theta=\cosh u\dd u$なので
\[
\int\sec^3\theta\dd\theta
=\int\sec\theta\cdot\sec^2\theta\dd\theta
=\int\cosh u\cdot\cosh u\dd u
=\int\cosh^2u\dd u
\]
である.三角置換で現れた$\sec\theta$の3乗の積分は,双曲線置換の$\cosh u$の2乗の積分と同じものだったのだ.三角置換では$\dd x$が$\sec^2\theta\dd\theta$となって次数が1つ上がり,部分積分が必要になった.双曲線置換では$\dd x=\cosh u\dd u$の次数が根号と同じなので,2倍角の公式だけで済んだ.

さらに,三角置換の答えに現れた$\sec\theta+\tan\theta$は$\cosh u+\sinh u=\e^u=t$である.3つの変数は
\[
x=\tan\theta=\sinh u,\qquad
t=\sec\theta+\tan\theta=\e^u
\]
で互いに移り合う.\chapterref{章:部分積分}で,\,$\{x+\sqrt{1+x^2}\}'$の中に元の$x+\sqrt{1+x^2}$が現れたのも,これが$t=\e^u$であり,\,$\frac{\dd t}{\dd u}=t$となるからである.

最後に答えを微分して確認しておく.\,$s=\sqrt{1+x^2}$と書けば$s'=\frac xs$であり,
\[
\frac{\dd}{\dd x}\log(x+s)
=\cfrac{1+\cfrac xs}{x+s}=\frac1s.
\]
したがって式\eqref{式:一と二乗の根号積分}の右辺の微分は
\[
\frac12\left(s+\frac{x^2}{s}+\frac1s\right)
=\frac12(s+s)=\sqrt{1+x^2}
\]
となる.3つの計算は同じ原始関数へ到達した.



%
%
\section*{章末問題}
\begin{enumerate}[label=\textbf{\thechapter.\arabic*},leftmargin=*]
\item\label{問:逆三角関数の枝}
$\arcsin(\sin\frac{4\pi}3)$,\,$\arccos(\cos\frac{5\pi}4)$,\,$\arctan(\tan(-\frac{2\pi}3))$を求めよ.また$y=\arcsin(x^2)$を$-1<x<1$で微分せよ.
\item\label{問:逆三角関数の積分練習}
$\int\frac1{x^2+2x+5}\dd x$と$\int_0^1\frac1{\sqrt{4-x^2}}\dd x$を求めよ.
\item\label{問:余接と正接の和}
定義から$(\cot x)'=-\csc^2x$を示せ.さらに$\frac1{\sin x\cos x}=\tan x+\cot x$を使って,\,$\int\frac1{\sin x\cos x}\dd x$を求め,答えを用いる区間の条件を述べよ.
\item\label{問:半角置換の区間分割}
半角置換を使って$\int_0^{2\pi}\frac1{2+\sin x}\dd x$を求めよ.\,$x=\pi$の前後を分け,置換後の無限の端点は極限として扱うこと.
\item\label{問:円の点とピタゴラス数}
円上の点$\left(\frac{12}{13},\frac5{13}\right)$に対応する傾き$t$を求めよ.\,$t=\frac nm$から式\eqref{式:ピタゴラス数}で作られる三整数を求め,共通因数で割れ.また$m,n$がともに奇数なら,三整数がすべて偶数になることを示せ.
\item\label{問:双曲線二倍角と積分}
指数関数の定義から$\sinh(2u)=2\sinh u\cosh u$を示せ.さらに$\int_0^{\log2}\cosh^2u\dd u$を求めよ.
\item\label{問:双曲線余弦の置換}
$x>1$で$\int\frac1{\sqrt{x^2-1}}\dd x$を,\,$x=\cosh u$によって求めよ.\,$u$の範囲を指定し,元の変数へ戻す対数の式も導け.
\item\label{問:双曲線置換と代数置換}
$\int_0^2\sqrt{4+x^2}\dd x$を,\,$x=2\sinh u$と$t=x+\sqrt{4+x^2}$の二通りの置換で求めよ.また二つの変数の間に$t=2\e^u$が成り立つことを確かめよ.
\end{enumerate}

\smallskip\noindent 解答は\pageref{解答:第11章}ページ.

%
